% Figure: RC circuit transient. The charging and discharging curves
% on a normalised axis (time in units of tau, voltage as fraction of V0).
\begin{figure}[htbp]
\centering
\begin{tikzpicture}[
  axis/.style={draw=engblack, -{Stealth}, thick},
  charge/.style={draw=engblue, very thick},
  discharge/.style={draw=engred, very thick},
  guide/.style={draw=enggray, dashed, thin},
  label/.style={font=\small},
  domain=0:5, samples=200,
]
% Axes
\draw[axis] (0,0) -- (10.5,0) node[anchor=west, label] {$t/\tau$};
\draw[axis] (0,0) -- (0,4.8)  node[anchor=south, label] {$v/V_{0}$};

% Tick marks on x-axis
\foreach \x/\m in {0/0, 2/1, 4/2, 6/3, 8/4, 10/5} {
  \draw (\x,-0.08) -- (\x,0.08);
  \node[font=\scriptsize, below] at (\x,-0.08) {\m};
}
% Tick marks on y-axis
\foreach \y/\h in {0/0, 1/0.25, 2/0.50, 3/0.75, 4/1.00} {
  \draw (-0.08,\y) -- (0.08,\y);
  \node[font=\scriptsize, left] at (-0.08,\y) {\h};
}

% Charging curve: y = 4 * (1 - exp(-x/2))   (4 units == 1.00)
\draw[charge] plot[domain=0:10, samples=200]
  ({\x}, {4*(1 - exp(-\x/2))});
% Discharging curve: y = 4 * exp(-x/2)
\draw[discharge] plot[domain=0:10, samples=200]
  ({\x}, {4*exp(-\x/2)});

% Reference horizontal lines at 0.63 and 0.37
\draw[guide] (0, 2.52) -- (10, 2.52);
\node[engblack, font=\scriptsize, anchor=west] at (10.05, 2.52) {$0.63$};
\draw[guide] (0, 1.48) -- (10, 1.48);
\node[engblack, font=\scriptsize, anchor=west] at (10.05, 1.48) {$0.37$};

% Vertical guide at one time constant (x=2)
\draw[guide] (2, 0) -- (2, 4);
\node[font=\scriptsize, below right, enggray] at (2, 0.5) {$t=\tau$};

% Labels on curves
\node[engblue, font=\small, anchor=south west] at (5.5, 3.6)
  {charging $v=V_{0}(1-e^{-t/\tau})$};
\node[engred, font=\small, anchor=north west] at (5.5, 0.4)
  {discharging $v=V_{0}e^{-t/\tau}$};

\end{tikzpicture}
\caption[RC circuit transient on the normalised time axis]{First-order
RC transients on the normalised axis $t/\tau$. The charging curve
(blue, from $0$ to $V_{0}$) and the discharging curve (red, from
$V_{0}$ to $0$) cross the $0.63 V_{0}$ and $0.37 V_{0}$ lines at one
time constant. The engineer who has internalised the percentage
calibration ($63\,\%$ at $1\tau$, $86\,\%$ at $2\tau$, $95\,\%$ at
$3\tau$, $99\,\%$ at $5\tau$) reads first-order settling time off
this picture by eye. Computed from $v=V_{0}(1-e^{-t/\tau})$ and
$v=V_{0}e^{-t/\tau}$.}
\label{fig:vol02:ch07:rc-transient}
\end{figure}
